% ===================================================================
%  TABELLA nr. 1 — La scola. --- L'institut. --- La sala da scola.
%
%  Ceci n'est PAS une transcription : c'est une traduction, faite en
%  2026, en rumantsch grischun. D'ou trois differences avec texto/io
%  et texto/fr, et trois seulement :
%
%   * pas de \nl ni de \cc — la traduction ne suit aucune ligne
%     imprimee, et les lignes de ce fichier ne sont que du confort ;
%   * pas de \begin{VUpage} — elle n'a ni feuillet ni folio, et la
%     page de lecture ne lui en affiche pas ;
%   * le gras se pose sur le TERME seul, ponctuation dehors.
%
%  Tout le reste est identique, et doit l'etre : les cles « %%K » sont
%  celles de texto/io, macro pour macro pour l'apparat, et les renvois
%  \textsuperscript{(n)} visent les MEMES objets de la planche — ce
%  sont eux qui ouvrent les gros plans.
%
%  LA SEIZIEME DES DIX-SEPT LANGUES IDISTES, et la premiere de toute
%  la serie dont la forme ecrite est PLUS JEUNE QUE LE LIVRET
%  TRADUIT. Le rumantsch grischun n'existe pas en 1926 : Heinrich
%  Schmid l'a compose en 1982 a partir des cinq idiomes ecrits des
%  vallees — sursilvan, sutsilvan, surmiran, puter, vallader. Les
%  quinze colonnes precedentes auraient pu, en droit, etre ecrites
%  l'annee du livret ; celle-ci non.
%
%  LE CHOIX DE L'IDIOME EST DONC UNE DECISION, ET LA VOICI. On prend
%  le rumantsch grischun et non le sursilvan, qui est le plus parle,
%  ni le vallader, qui a la plus vieille litterature, parce que c'est
%  la seule forme qui vaut pour tout le domaine et parce que c'est la
%  seule qui soit ecrite pour l'ECOLE : la langue commune a ete faite,
%  d'abord, pour qu'on puisse imprimer un manuel unique. Le tableau 1
%  est une salle de classe : la colonne s'ouvre exactement la ou le
%  rumantsch grischun a commence.
%
%  ET LA LANGUE LA PLUS JEUNE PORTE LES MOTS LES PLUS VIEUX. Trois
%  mots de cette page le montrent :
%    « cudesch »  le livre, du latin CODEX — l'italien dit « libro »,
%                 le francais « livre », tous deux de LIBER ;
%    « pled »     le mot, du latin PLACITUM — l'italien dit
%                 « parola », le francais « mot », ni l'un ni l'autre
%                 n'a garde PLACITUM pour cela ;
%    « encaust »  l'encre, du grec ENCAUSTUM par le latin, la ou le
%                 francais a use le meme mot jusqu'a « encre » et ou
%                 l'italien a pris « inchiostro ».
%  UNE NORME DE 1982 N'EST PAS UN VOCABULAIRE DE 1982 : ce que Schmid
%  a fixe, c'est l'orthographe commune, pas le fonds, et le fonds est
%  d'avant que l'italien et le francais existent.
%
%  ON TRADUIT SUR L'IDO ET L'ON CONTROLE SUR texto/fr, texto/ca et
%  texto/oc. On ne recopie aucune des trois : « aula » se dit « sala
%  da scola », « professor » se dit « magister », « llibre » se dit
%  « cudesch », « paraula » se dit « pled ».
%
%  MEME ORDRE DES MOTS QUE DANS LES AUTRES ROMANES, donc meme ordre
%  des renvois : le complement du nom se postpose, et l'ido s'y ecrit
%  sans qu'un seul renvoi ait a bouger.
%
%  LES NOMS DE PERSONNE SONT ROMANCHES : Ioannes est Gion, Iakobus
%  est Jachen, Georgius est Gieri, Franciskus est Franzestg. Les noms
%  latins de la note (*) ne bougent pas : c'est d'eux que la note
%  parle.
% ===================================================================

%%K t01-apar-0 apar
\VUsaut{2.0mm}
\VUcentre{12.6pt}{120}{LIBRET EXPLICATIV}
\VUsaut{3.2mm}
\VUcentre{8.4pt}{160}{DA}
\VUsaut{2.6mm}
\VUtitre{15.6pt}{0}{las tabellas auxiliaras Delmas}
\VUsaut{3.4mm}
\VUornamento{9pt}
\VUsaut{5.0mm}
\VUcentre{13.2pt}{90}{EMPRIMA SERIA}
\VUsaut{3.0mm}
\VUfilet{16mm}
\VUsaut{5.6mm}

%%K t01-apar-1 apar
\VUtitre{15.0pt}{60}{TABELLA nr. 1}
\VUsaut{1.4mm}
\VUtitre{11.4pt}{0}{La scola. --- L'institut. --- La sala da scola.}
\VUsaut{2.2mm}
\VUfilet{20mm}
\VUsaut{6.4mm}

%%K t01-tit-1 sub
\VUpk{10.2pt}{40}{Descripziun generala.}
\VUsaut{3.0mm}

%%K t01-01-1 p
1. --- Questa tabella represchenta ina \VUgras{sala da scola} en in
\VUgras{institut}. En questa sala da scola i ha otg \VUgras{maisas}
\textsuperscript{(18)} ed otg \VUgras{bancs} \textsuperscript{(32)}. Sin mintga
banc sesan trais scolars. Il \VUgras{magister} \textsuperscript{(7)} sesa sin ina
\VUgras{sutga} \textsuperscript{(8)} a sia \VUgras{catedra} \textsuperscript{(6)}.

%%K t01-02-1 p
2. --- A sanestra da la catedra sa chatta in \VUgras{armari} cun portas da vaider
\textsuperscript{(15)}. A dretga è la \VUgras{tavla naira} \textsuperscript{(1)} cun
il \VUgras{sfruschader} \textsuperscript{(2)} e la \VUgras{crida}
\textsuperscript{(3)}.

%%K t01-02-2 p
Sur la catedra pendan \VUgras{tavlas da mir} \textsuperscript{(9, 11, 12)} ed ina
\VUgras{charta geografica} \textsuperscript{(10)}.

%%K t01-03-1 p
3. --- Betg tut ils scolars sesan a lur \VUgras{plaz} \textsuperscript{(17)}. Gion
\textsuperscript{(4)} stat sper la tavla; el tegna il sfruschader e stizza las
\VUgras{figuras geometricas}.

%%K t01-03-2 p
Peider è dasperas la catedra; el rimna ils \VUgras{exercizis scrits}
\textsuperscript{(5)} ed als porta al magister.

%%K t01-04-1 p
4. --- \VUgras{Quest} magister è il magister da \VUgras{linguas vivas}. El
instruescha als scolars da discurrer, leger e scriver linguas estras
\VUgras{(tudestg, englais, spagnol, talian, russ} u \VUgras{franzos)}.

%%K t01-05-1 p
5. --- La sala da scola è fitg vasta e fitg clera. En il fund i ha ina largia
\VUgras{fanestra}, dus \VUgras{vasistas} \textsuperscript{(78)} ed ina \VUgras{porta
da vaider} per l'arieschar e l'illuminar.

%%K t01-05-2 p
Questa sala da scola n'è betg fraida. En il mez i ha, per la stgaudar, in fitg bun
\VUgras{fuorn} \textsuperscript{(46)} cun ses \VUgras{tub} \textsuperscript{(45)}.

%%K t01-05-3 p
Vi da la paraid èn fixadas \VUgras{hoccas da vestgadira} \textsuperscript{(58)}; vi
da quellas pendan las bricolas ed ils mantels dals scolars.

%%K t01-tit-2 sub
\VUsaut{5.2mm}
\VUpk{10.2pt}{40}{La curt da recreaziun.}
\VUsaut{3.6mm}

%%K t01-06-1 p
6. --- L'\VUgras{ura} \textsuperscript{(63)} mussa nov uras e tschintg minutas.

%%K t01-06-2 p
La \VUgras{recreaziun} \textsuperscript{(76)} è gist finida e la lecziun cumenza
danovamain. La recreaziun ha lieu en la \VUgras{curt} \textsuperscript{(75)}. Nus
vesain intgins scolars che giogan. In \VUgras{sutmagister} \textsuperscript{(74)}
als survegliescha.

%%K t01-07-1 p
7. --- En la curt vesain nus dapli persunas: l'\VUgras{inspectur}
\textsuperscript{(73)}, il \VUgras{directur} \textsuperscript{(71)} ed il
\VUgras{censur} \textsuperscript{(72)}.

%%K t01-07-2 p
L'inspectur inspecziunescha las scolas, il directur dirigescha l'institut, il
censur survegliescha ils studis.

%%K t01-07-3 p
La quarta persuna è il \VUgras{portier} \textsuperscript{(77)}. El porta sut il
bratsch in register per inscriver ils absents.

%%K t01-08-1 p
8. --- En il fund da la curt sa chattan ils \VUgras{apparats da gimnastica}
\textsuperscript{(70)} e l'\VUgras{ovratori} da \VUgras{lavurs manualas}
\textsuperscript{(69)}.

%%K t01-08-2 p
A dretga da la tabella cumenzain nus a vesair las \VUgras{salas} da \VUgras{musica}
e da \VUgras{dissegn}.

%%K t01-noto-1 noto
\VUnotes{143.2mm}{%
(*) Per ils \textit{nums da batten} vegn cusseglià da als transcriver era en la
furma latina. Quai è l'unic med per evitar variantas senza dumber. Uschiglio, co
eleger tranter \textit{Giovanni, Jean,} \textit{Johann, Jan, Hans, John, Ivan,}
etc.? Stuain nus crear in dicziunari spezial per ils nums da batten? Nus
prendain dal latin lur nominativ e din: \VUgras{Ioannes, Ioanna; Iulius, Iulia;
Iakobus; Andreas; Lukas.} (Grammatica cumpletta).%
}

%%K t01-tit-3 sub
\VUpk{10.2pt}{40}{Descripziun detagliada.}
\VUsaut{2.4mm}

%%K t01-tit-4 sub
\VUpk{10.2pt}{40}{Ils buns scolars.}
\VUsaut{3.0mm}

%%K t01-09-1 p
9. --- Ils scolars da l'emprim banc a dretga da la catedra èn \VUgras{Enric} (*),
\VUgras{Stefan} e \VUgras{Carli}.

%%K t01-09-2 p
Enric è fitg attentiv e lavurus; el asculta il magister. Sin sia maisa ha el in
\VUgras{quadern} \textsuperscript{(28)}, in \VUgras{cudesch} \textsuperscript{(29)},
in \VUgras{vocabulari} \textsuperscript{(30)} ed ina \VUgras{chaschetta da penns}
\textsuperscript{(31)}. Ses cudesch ha bleras \VUgras{paginas}; mintga chapitel
cuntegna pliras \VUgras{paragrafs} e mintga \VUgras{lingia} blers \VUgras{pleds}.

%%K t01-09-4 p
Ses vischin Stefan \textsuperscript{(24)} è fervent. El nota en ses quadern la
lecziun per la proxima giada. El ha ina bella \VUgras{scrittira}. Ses cudesch è
serrà. Nus vesain mo la \VUgras{cuvertura} \textsuperscript{(27)}. Sper el sa chatta
ses \VUgras{cumpass} \textsuperscript{(26)}.

%%K t01-09-5 p
Carli \textsuperscript{(23)} tegna ses cudesch en maun; el al avra per repeter sia
lecziun. Sco mintga scolar ha el in \VUgras{encausteri} \textsuperscript{(25)}
avant el. El è obedient.

%%K t01-10-1 p
10. --- A la maisa dasperas èn \VUgras{Luis, Antoni} e \VUgras{Pauli}. Luis
recitescha sia \VUgras{lecziun} \textsuperscript{(22)} e responda a las
\VUgras{dumondas} dal magister. Antoni \textsuperscript{(21)} è fitg distratg;
mintgatant al chastiescha il magister schizunt. Pauli \textsuperscript{(20)} è in
bun \VUgras{cumpogn}, amiaivel e servetsch. El è fitg attentiv.

%%K t01-tit-5 sub
\VUsaut{4.2mm}
\VUpk{10.2pt}{40}{Ils mals scolars.}
\VUsaut{3.2mm}

%%K t01-11-1 p
11. --- Suenter Enric sa chatta \VUgras{Gieri} \textsuperscript{(38)}. El n'è betg
in nausch mat, ma el è \VUgras{tschintschader}, distratg e malquiet. Sia
\VUgras{superficialitad} e sia \VUgras{malobedienza} chaschunan \VUgras{disturbi}
durant la lecziun, ed el merita savens in sever \VUgras{avertiment}. Ses
\VUgras{quadern da sbozz} \textsuperscript{(35)} è nausch, el l'emplenescha da
\VUgras{maglias d'encaust} cun sia \VUgras{penna} \textsuperscript{(34)}, e perquai
duvra el adina sia \VUgras{gumma da stizzar} \textsuperscript{(36)}; sia
\VUgras{chartatscha} \textsuperscript{(33)} è squarschada, perquai ch'el è fitg
negligent. Pertge porta el ses \VUgras{portamina} \textsuperscript{(37)} en la sala
da scola da linguas vivas?

%%K t01-11-2 p
Ses vischin Edmund \textsuperscript{(39)} n'al ha betg gugent. El è propi in pau
pigher, ma el è taschaivel e sa tegna quiet. El na tschintscha betg; mo, empè
d'ascultar, preferescha el da leger las \VUgras{istorgias} da ses \VUgras{cudesch
da lectura}.

%%K t01-11-3 p
Il terz scolar da quest banc è \VUgras{Luis} \textsuperscript{(40)}. El è diligent.
El scriva sin in \VUgras{fegl da palpiri} alv. Avant el ha el mess ses
\VUgras{palpiri sugant} \textsuperscript{(41)}.

%%K t01-12-1 p
12. --- Ils scolars dal segund banc, a sanestra dal magister, èn fitg giuvens.
L'emprim, il pitschen \VUgras{Emil}, na sa anc betg scriver bain; el porta cun el
ses \VUgras{tavlin d'ardesia} \textsuperscript{(42)}, ses \VUgras{stilus} e sia
\VUgras{spugna} \textsuperscript{(43)}.

%%K t01-12-2 p
Sper el èn \VUgras{August} e \VUgras{Bernard} \textsuperscript{(44)}. Quest ultim
lavura bain, ma sia \VUgras{vestgadira} è fitg negligida.

%%K t01-13-1 p
13. --- Suenter els sa chattan trais \VUgras{pigers. Albert} n'emprenda betg sias
lecziuns. \VUgras{Jachen} \textsuperscript{(47)} dorma empè d'ascultar. Sia
\VUgras{pigrezza} al procura \VUgras{reproschas} e schizunt \VUgras{chastis}.
Finalmain è \VUgras{Cristian} \textsuperscript{(48)} memia sfrenà. El
\VUgras{sa cumporta} mal e na fa nagin sforz per sa meglierar.

%%K t01-14-1 p
14. --- Ses vischins, percunter, sa cumportan bain. \VUgras{Ernest}
\textsuperscript{(51)} ha gugent l'\VUgras{urden} e la \VUgras{regularitad}; el
duvra adina sia \VUgras{regla} \textsuperscript{(50)} e ses \VUgras{crajun}
\textsuperscript{(49)}. El è \VUgras{extern}.

%%K t01-14-2 p
\VUgras{Franzestg} \textsuperscript{(52)} è \VUgras{intern}; el porta
l'\VUgras{uniforma} da l'institut; el maglia e dorma là. El è in bun
\VUgras{cumpogn}.

%%K t01-14-3 p
Maurizi taglia ses \VUgras{crajun} cun ses \VUgras{curtelin}
\textsuperscript{(56)}; el è fitg premurus.

%%K t01-14-4 p
El porta cun el tut ses cudeschs, sia \VUgras{grammatica} \textsuperscript{(55)},
ses \VUgras{cudeschet da notas} \textsuperscript{(54)} e schizunt in
\VUgras{sutmaun} \textsuperscript{(53)}. Ses \VUgras{satg da scola}
\textsuperscript{(57)} è per terra suenter el.

%%K t01-14-5 p
Las duas maisas en il fund da la sala da scola èn occupadas da \VUgras{buns
scolars} \textsuperscript{(19)}.

%%K t01-tit-6 sub
\VUsaut{4.6mm}
\VUpk{10.2pt}{40}{Dissegn e musica.}
\VUsaut{3.4mm}

%%K t01-15-1 p
15. --- En la \VUgras{sala da dissegn} i ha bels \VUgras{models} pendids vi da la
paraid ed ina \VUgras{lampa} \textsuperscript{(62)} cun \VUgras{para-glisch},
suspendida al plafond. Il \VUgras{magister} \textsuperscript{(60)} è fitg pazient;
el curregia ils \VUgras{dissegns} \textsuperscript{(61)} dals scolars ed als
instruescha da dissegnar in \VUgras{bust} \textsuperscript{(59)} tenor la natira.

%%K t01-16-1 p
16. --- La \VUgras{sala da musica} para fitg interessanta. Qua emprenda ins da
leger la \VUgras{musica}, da solfegiar las \VUgras{notas da la gamma}
\textsuperscript{(67)} scrittas sin il \VUgras{sistem da lingias} (do, re, mi, fa,
sol, la, si\,=\,C. D. E. F. G. A. B.), da renconuscher las \VUgras{clavs} e da
chantar bellas \VUgras{melodias}. Ins metta las partituras sin il \VUgras{pult}
\textsuperscript{(65)}. In dals scolars \VUgras{suna il piano}
\textsuperscript{(64)} ed il \VUgras{magister da musica} \textsuperscript{(66)}
\VUgras{bat la mesira} \textsuperscript{(68)}.

%%K t01-17-1 p
17. --- Nus cumenzain a vesair in chantun da l'\VUgras{ovratori} da \VUgras{lavurs
manualas} \textsuperscript{(69)}, nua ch'ils mats pon emprender da plajar il lain e
da limar il fier. Quai n'exista betg en tut ils instituts.

%%K t01-tit-7 sub
\VUsaut{5.0mm}
\VUpk{10.2pt}{40}{La sala da scienzas.}
\VUsaut{3.6mm}

%%K t01-18-1 p
18. --- Il magister da matematica instruescha als scolars la \VUgras{geometria,
l'aritmetica, il calcul} e il \VUgras{sistem metric}. Nus vesain sin la tabella las
figuras geometricas principalas: in \VUgras{circul} \textsuperscript{(a)}, in
\VUgras{quadrat} \textsuperscript{(b)}, in \VUgras{triangul} \textsuperscript{(c)},
ina \VUgras{lingia} \textsuperscript{(d)}.

%%K t01-18-2 p
Nus vesain era ina \VUgras{operaziun}; quai n'è betg ina \VUgras{adiziun}, ni ina
\VUgras{sutracziun}, ni ina \VUgras{multiplicaziun}: quai è ina \VUgras{divisiun}
\textsuperscript{(e)}.

%%K t01-19-1 p
19. --- Sin la tavla dal sistem metric distinguin nus: in \VUgras{meter}
\textsuperscript{(a)}, ina \VUgras{balanza} \textsuperscript{(b)} e sias
\VUgras{peisas}, in \VUgras{liter} \textsuperscript{(e)}, in \VUgras{decaliter}
\textsuperscript{(f)}, in \VUgras{deciliter} ed in \VUgras{meter cubic}
\textsuperscript{(d)}.

%%K t01-19-2 p
Ils scolars studegian era las \VUgras{scienzas natiralas} \textsuperscript{(11)}
--- zoologia \textsuperscript{(a)}, botanica \textsuperscript{(b)}, geologia
\textsuperscript{(c)} --- e l'\VUgras{istorgia} \textsuperscript{(12)}.

%%K t01-19-3 p
En l'armari sa chattan ils apparats da \VUgras{fisica} \textsuperscript{(15)} e da
\VUgras{chemia} \textsuperscript{(16)}.

%%K t01-19-4 p
Sur l'armari sa chattan: ina \VUgras{sfera} \textsuperscript{(14)}, in
\VUgras{cono} ed in \VUgras{globus terrestric} \textsuperscript{(13)}.

%%K t01-19-5 p
Sur las tavlas penda ina \VUgras{charta} che represchenta las tschintg parts dal
mund: l'\VUgras{Europa} \textsuperscript{(g)}, l'\VUgras{Asia}
\textsuperscript{(e)}, l'\VUgras{Africa} \textsuperscript{(f)},
l'\VUgras{America} \textsuperscript{(ab)} e l'\VUgras{Oceania}
\textsuperscript{(h)}, ed ils dus gronds oceans, il \VUgras{Pacific}
\textsuperscript{(c)} e l'\VUgras{Atlantic} \textsuperscript{(d)}.
\VUsaut{6.0mm}
\VUfilet{22mm}
